\documentclass[11pt,a4paper]{article}
\usepackage[utf8]{inputenc}
\usepackage[T1]{fontenc}
\usepackage[brazilian]{babel}
\usepackage{geometry}
\usepackage{enumitem}
\usepackage{booktabs}
\usepackage{amsmath}
\usepackage{graphicx}
\graphicspath{{../figures_global/}}
\usepackage{fancyhdr}
\usepackage[backend=biber, style=numeric, sorting=none]{biblatex}
\usepackage{hyperref}
\addbibresource{referencias.bib}

\geometry{margin=2.5cm, headheight=32pt}
\setlength{\parindent}{0pt}
\setlength{\parskip}{0.5em}

\pagestyle{fancy}
\fancyhf{}
\fancyhead[L]{\small\textbf{Faculdade de Ciências e Tecnologias em Engenharia}}
\fancyhead[R]{\includegraphics[height=0.85cm]{../figures_global/unb.png}}
\renewcommand{\headrulewidth}{0.4pt}
\fancyfoot[C]{\small Princípios de Comunicação para Engenharia --- Ementa}
\fancyfoot[R]{\small \thepage}
\renewcommand{\footrulewidth}{0.4pt}
\fancypagestyle{plain}{%
  \fancyhf{}%
  \fancyhead[L]{\small\textbf{Faculdade de Ciências e Tecnologias em Engenharia}}%
  \fancyhead[R]{\includegraphics[height=0.85cm]{../figures_global/unb.png}}%
  \renewcommand{\headrulewidth}{0.4pt}%
  \fancyfoot[C]{\small Princípios de Comunicação para Engenharia --- Ementa}%
  \fancyfoot[R]{\small \thepage}%
  \renewcommand{\footrulewidth}{0.4pt}%
}

\title{\textbf{Ementa}\\[0.3em]\large Princípios de Comunicação para Engenharia}
\author{Faculdade de Ciências e Tecnologias em Engenharia --- FCTE/UnB}
\date{}

\begin{document}

\thispagestyle{plain}
\maketitle
\vspace{-1.0em}

\section{Identificação}

\begin{center}
\small
\begin{tabular}{@{}p{0.32\textwidth} p{0.60\textwidth}@{}}
\toprule
\textbf{Componente curricular} & Princípios de Comunicação para Engenharia \\
\textbf{Natureza} & Obrigatória \\
\textbf{Créditos} & 4 créditos --- 3 teóricos e 1 de extensão \\
\textbf{Carga horária total} & 60 horas \\
\textbf{Carga horária teórica} & 45 horas \\
\textbf{Carga horária de extensão} & 15 horas \\
\textbf{Modalidade} & Presencial \\
\textbf{Pré-requisitos} & Sinais e Sistemas; Teoria de Circuitos Eletrônicos I \\
\bottomrule
\end{tabular}
\end{center}

A carga horária de extensão atende à curricularização da extensão prevista na Resolução CNE/CES 07/2018 e nas normas da UnB, correspondendo a 25\% da carga horária total do componente.

\section{Ementa}

Introdução aos sistemas de comunicação e nivelamento em análise de sinais. Diagrama de blocos de um sistema de comunicação. Classificação de sinais. Revisão de Série e Transformada de Fourier. Densidade espectral de energia e de potência. Sinais passa-banda e representação em fase e quadratura. Modulação analógica em amplitude e suas variantes. Modulação angular em frequência e em fase. Malha de captura de fase. Ruído em sistemas de comunicação analógicos e relação sinal-ruído em demoduladores. Amostragem, quantização e modulação por código de pulso. Codificação de linha e interferência entre símbolos. Desempenho de enlaces digitais e taxa de erro de bit. Noções de protocolos e arquiteturas de redes. Projeto de extensão aplicado a demanda real da comunidade externa, com desenvolvimento e entrega de um produto de comunicação.

\section{Objetivos}

\subsection*{Objetivo geral}

Capacitar o estudante a analisar, projetar e avaliar os blocos fundamentais de um sistema de comunicação, articulando a formação teórica com a solução de um problema real por meio de projeto de extensão.

\subsection*{Objetivos específicos}

\begin{itemize}[noitemsep]
  \item Compreender a análise espectral de sinais e sua aplicação em sistemas de comunicação.
  \item Dominar as técnicas de modulação analógica em amplitude, frequência e fase.
  \item Analisar o efeito do ruído em sistemas analógicos e digitais de comunicação.
  \item Entender os processos de amostragem, quantização e modulação por código de pulso.
  \item Avaliar o desempenho de enlaces digitais em termos de taxa de erro de bit.
  \item Obter noções básicas de protocolos e arquiteturas de redes de comunicação.
  \item Aplicar os conceitos da disciplina no desenvolvimento de um produto de comunicação para uma demanda concreta da comunidade externa.
\end{itemize}

\section{Conteúdo Programático}

\begin{center}
\small
\begin{tabular}{|p{0.30\textwidth}|p{0.30\textwidth}|p{0.30\textwidth}|}
\hline
\textbf{Unidade 1: Introdução e Nivelamento} & \textbf{Unidade 2: Modulações Analógicas} & \textbf{Unidade 3: ADC e Comunicação Digital} \\
\hline
\parbox[t]{0.28\textwidth}{%
  \vspace{0.2em}
  \begin{itemize}[noitemsep, leftmargin=*, nosep, font=\small]
    \item Sistemas de comunicação e diagrama de blocos
    \item Classificação de sinais
    \item Revisão de Série e Transformada de Fourier
    \item Densidade espectral de energia e de potência
    \item Sinais passa-banda
  \end{itemize}
  \vspace{0.2em}
} &
\parbox[t]{0.28\textwidth}{%
  \vspace{0.2em}
  \begin{itemize}[noitemsep, leftmargin=*, nosep, font=\small]
    \item Modulação em amplitude AM
    \item Variantes DSB, DSB+C, SSB, VSB
    \item Modulação em frequência FM
    \item Modulação em fase PM
    \item Malha de captura de fase PLL
    \item Ruído em sistemas analógicos
    \item Relação sinal-ruído em demoduladores
  \end{itemize}
  \vspace{0.2em}
} &
\parbox[t]{0.28\textwidth}{%
  \vspace{0.2em}
  \begin{itemize}[noitemsep, leftmargin=*, nosep, font=\small]
    \item Amostragem
    \item Quantização
    \item PCM linear e não linear
    \item Codificação de linha
    \item Interferência entre símbolos
    \item Taxa de erro de bit
    \item Noções de protocolos
  \end{itemize}
  \vspace{0.2em}
} \\
\hline
\end{tabular}
\end{center}

\section{Componente de Extensão}

A componente de extensão corresponde a 1 crédito e organiza-se como um projeto desenvolvido ao longo do semestre, no qual os estudantes aplicam o conteúdo da disciplina para atender a uma demanda real de um ator externo à universidade. O projeto é entregue ao final do semestre e deve resultar em um produto ou solução de comunicação com valor concreto para o parceiro.

\subsection*{Formas de interação com a comunidade externa}

\begin{itemize}[noitemsep]
  \item Parceria com empresas do setor de eletrônica, telecomunicações ou correlatos, a partir de um problema por elas apresentado.
  \item Apoio técnico a equipes de competição da universidade, como foguetemodelismo, satélites universitários, veículos autônomos e afins, no desenvolvimento de enlaces de telemetria e comunicação.
  \item Atendimento a demandas de instituições públicas, associações comunitárias ou grupos de radioamadorismo.
\end{itemize}

\subsection*{Exemplos de projetos de extensão}

\begin{itemize}[noitemsep]
  \item Enlace de telemetria em radiofrequência para foguete ou minissatélite de equipe de competição.
  \item Sistema de recepção e decodificação de sinais de satélites meteorológicos para uso didático e comunitário.
  \item Módulo de transmissão de dados de baixo custo para monitoramento ambiental em parceria com órgão externo.
  \item Prova de conceito de enlace de comunicação demandada por empresa parceira, com caracterização de desempenho.
\end{itemize}

\subsection*{Produtos e registro da extensão}

Cada grupo entrega um relatório técnico do projeto, um protótipo ou demonstração funcional e uma apresentação final aberta ao parceiro externo. As atividades e a interação com a comunidade são registradas conforme as normas de extensão da UnB, incluindo o cadastro da ação e o registro da carga horária dos estudantes.

\section{Metodologia}

As aulas teóricas combinam exposição dialogada, resolução de problemas e demonstrações computacionais. A componente de extensão é conduzida por aprendizagem baseada em projetos, com acompanhamento em encontros de orientação, definição de escopo junto ao parceiro externo e entregas parciais ao longo do semestre. O material da disciplina é disponibilizado pelo ambiente virtual institucional.

\section{Avaliação}

A nota final combina a componente teórica e a componente de extensão.

\subsection*{Componente teórica}

A nota teórica $N_T$ é a média ponderada de três provas:
\[
N_T = \frac{P_1 + P_2 + 2\,P_3}{4}
\]

\subsection*{Componente de extensão}

A nota de extensão $N_E$ resulta do projeto desenvolvido junto ao parceiro externo, avaliado por entregas parciais, relatório técnico final, protótipo ou demonstração e apresentação:
\[
N_E = 0{,}3\,\text{Parciais} + 0{,}4\,\text{Relatório} + 0{,}3\,\text{Apresentação}
\]

\subsection*{Nota final e aprovação}

\[
N_F = 0{,}75\,N_T + 0{,}25\,N_E
\]

\begin{itemize}[noitemsep]
  \item Exige-se aprovação simultânea na teoria e na extensão, com $N_T \geq 5$ e $N_E \geq 5$.
  \item A aprovação final é obtida com $N_F \geq 5$.
  \item É obrigatória a frequência mínima de 75\% nas atividades teóricas e de extensão.
\end{itemize}

\section{Referências}

\nocite{lathi2019modern, haykin2009communication, proakis2007fundamentals, couch2013digital, carlson2002communication, ziemer2015principles}
\printbibliography[heading=none]

\end{document}
